\scanpage{10}
% archive foliation: 9
\deleted{un sous-groupe cyclique d'ordre 4, noté $\mathfrak{D}^+$,
et $S$ et $A$ sont également des torseurs sous \unsure{$\mathfrak{D}$}
$\mathfrak{D}^+$.}
groupe d'ordre 2 \deleted{central} qui est le centre, et dont l'élément
$\ne1$ est \uline{l'antipodisme} $\underline{a}$.

Il y a dans $\mathfrak{D}$ exactement deux éléments \inserted{$u,u'$}
d'ordre 4, inverses l'un de l'autre, engendrant l'unique sous-groupe cyclique
d'ordre 4, noté $\mathfrak{D}^+$ --- et $S$ et $A$ sont des torseurs sous ce
même groupe. L'antipodisme est l'unique élément d'ordre 2 de $\mathfrak{D}^+$.
\margin{c'est lui $\Gamma$}
\deleted{Le quotient $\mathfrak{D}/\mathfrak{D}^+$ opère sur}
\[
\deleted{\unsure{\mathfrak{D}_{\mathrm{ab}}^+\simeq
\mathfrak{D}/\{1,\underline{a}\}}}
\]

Le quotient $\mathfrak{D}/V_S\simeq\deleted{\unsure{\mathbb{F}_2}}\mathbb{Z}/2\mathbb{Z}$
opère sur $V_S$ par l'identité sur
$\{1,\underline{a}\}\simeq\mathbb{F}_2\subset V_S\simeq\mathbb{F}_2^C$,
et par la transposition sur $V_S\setminus\mathbb{F}_2\simeq C$,
et symétriquement pour $\mathfrak{D}/V_A\simeq\mathbb{Z}/2\mathbb{Z}$ opérant sur
$V_A\simeq\mathbb{F}_2^D$, \deleted{\unsure{respectivement la transpo}}
via l'opération par transposition sur $D$.

Le quotient $\mathfrak{D}/\mathfrak{D}^+\simeq\mathbb{Z}/2\mathbb{Z}$ opère
sur $\mathfrak{D}^+$ par l'autom.
\deleted{\unsure{de ce groupe}} $x\mapsto x^{-1}$,
\deleted{\unsure{en permutant}}\inserted{en particulier}
$u\mapsto u'=u^{-1}$ i.e. l'opération induit la transposition sur l'ens
$\underline{\omega}_Q$ des deux générateurs de
\deleted{\unsure{$\mathbb{Z}$}} $\mathfrak{D}^+$.
\margin{$\mathfrak{D}^+\simeq
\mathbb{Z}/4\mathbb{Z}\wedge_{\{\pm1\}}\underline{\omega}_Q$}

On a
\[
\mathfrak{D}_{\mathrm{ab}}^+
 \simeq \mathfrak{D}/\{1,\underline{a}\},
\]
espace vect. de rg 2 sur $\mathbb{F}_2$.
Les trois sous-groupes $V_S,\mathfrak{D}^+,V_A$ sont les images inverses des
trois sous-groupes d'ordre 2 de
$\mathfrak{D}_{\mathrm{ab}}^+$,
donc les trois seuls sous-groupes d'indice 2 de
$\mathfrak{D}_{\mathrm{ab}}^+$.
On a une base canonique de
$\mathfrak{D}_{\mathrm{ab}}^+$
\[
\mathfrak{D}_{\mathrm{ab}}^+
 \simeq\mathbb{F}_2\times\mathbb{F}_2
\]
définie par la condition que

\margin{$\mathfrak{D}$ opère sur $C$ et sur $D$ par autom. intérieurs,
donc sur $C\cup D$ en fixant $C,D$, d'où
$\mathfrak{D}\longrightarrow
\unsure{\mathcal{G}_C|\mathfrak{S}_C}\times\unsure{\mathcal{G}_D|\mathfrak{S}_D}
\simeq\mathbb{F}_2+\mathbb{F}_2$ est surjectif, et a un noyau
$\{1,\underline{a}\}$, et définit l'isom. can.
$\mathfrak{D}/\{1,\underline{a}\}=\mathbb{F}_2\times\mathbb{F}_2$.
\deleted{\uline{NB}
$\unsure{\mathfrak{D}_{\mathrm{ab}}^+\simeq
V_S\times V_A/\mathbb{F}_2\simeq
\mathbb{F}_2^C\times\mathbb{F}_2^D}$}}
\margin{\unsure{et $D$ opère sur $C$ comme $D$ opère sur $C$}
(i.e. $c\in C$, $d\in D\Longrightarrow[c,d]=\underline{a}$)}
